\section{Motivation}


\revise{Software systems are increasingly used to manage real-world
assets of significant value. They hold currencies and the digital
wallets that store them, underpin the operations of banks and other
financial institutions, and keep much of the sensitive personal
information that was once recorded on paper. 
This trend is large and
accelerating: the IMF's 2025 Financial 
Access Survey reports that
digital transactions, including mobile 
money and mobile and internet
banking, have surged in emerging and 
developing economies, rising
from 55 transactions per adult in 2017 to 251 per adult by
2024~\cite{imf2025fas}. As more value 
moves into software, a single
defect or compromise can translate directly into 
real-world harm.}

\revise{Not only has software itself changed, but so 
has the way
software is built. Despite well-known
limitations of their black-box nature,
it is by now an
undeniable fact that developers 
routinely rely on AI-assisted coding
tools to help implement software. These 
tools dramatically lower the
barrier to development: as long as a 
need can be described in natural
language, the tool can try its best to 
help produce a 
corresponding implementation.
This opens software development to a far wider audience, including
people with little formal programming 
background. Tools such as Claude
Code~\cite{anthropic_claude_code}, Codex~\cite{openai_codex_2026}, and
OpenCode~\cite{opencode_2026} are now widely 
used in practice. }

\revise{Not only has the way software is built changed, but so has the
question of who, or what, uses it.
Traditionally, software was
designed with specific users in mind: when 
building an application, a
developer could reasonably assume it would 
be operated by humans, for
instance tapping on an iPhone screen. 
This assumption no longer holds.
As soon as software exposes 
programmable interfaces, those interfaces
can be invoked by autonomous 
agents in arbitrary ways that the
developers may never have 
anticipated at design time.}

\revise{The above two facts, despite the convenience and benefits they
bring, open new risks compared with the traditional workflow of
developing software. Figure~\ref{fig:workflow} further clarifies the
new risks introduced on top of the existing software development
lifecycle. First, in the software development stage, not only do humans
make mistakes in the limited number of lines they write themselves, but
errors can also arise in the thousands of lines of code generated by
AI-assisted coding tools. Such a large volume of code makes it
difficult even for expert developers to review, let alone the new
non-technical programmers who have not received sufficient software
training. Second, in the software deployment stage, when monitoring
users' activity, traditional metrics such as click rates and
day--night usage patterns no longer apply. Human users may only use
social media when they are awake, which follows a typical biological
cycle, but autonomous agents can access the software in unprecedented
ways. This new usage, often deviating from developers' original intent,
can pose serious risks to software that stores critical real-world
value, such as digital currency, credentials, and classified
blueprints.}

\begin{figure}[t]
  \centering
  \includegraphics[width=\textwidth]{figures/workflow.pdf}
  \caption{\revise{The overall workflow for developing
  emerging software systems.
  In \emph{Stage~1 (development)}, software is written by
  human developers together with AI-assisted coding tools, which
  introduce new risks into the code itself. After review, testing,
  auditing, and verification, the software moves to \emph{Stage~2
  (deployment)}, where it is invoked not only by human users but
  also by autonomous agents, introducing further risks.}}
  \label{fig:workflow}
\end{figure}

\revise{Because of these shifts, we are already witnessing rapidly
growing security threats to software. Cyberattacks are not a new
problem. The IMF reports that the financial sector has lost \$12
billion over the past two decades to more than 20{,}000 such
incidents~\cite{financemagnates2024imf}. However, the two changes
above sharply amplify the threat. On the development side,
AI-generated code has already caused direct financial loss. In early
2026, a faulty price-oracle update in a pull request co-authored by an
AI coding model mispriced collateral on the Moonwell lending protocol,
leading to roughly \$1.78 million in losses~\cite{moonwell2026exploit}.
On the usage side, autonomous agents have themselves been turned into
attack tools. A state-sponsored group weaponized an AI coding agent to
run a largely autonomous cyber-espionage campaign against around
thirty organizations before the provider detected the activity and
disabled the accounts~\cite{anthropic2025espionage}. These attacks
compound an already worsening landscape. In 2025, a record year for
cryptocurrency theft, attackers stole more than \$3 billion in digital
assets~\cite{chainalysis2025crypto}.}

\revise{It is therefore increasingly important to protect what we call
\emph{emerging software systems}. Informally, an emerging software
system is durable, high-stakes software whose purpose is to protect
and operate on critical assets such as digital wallets, secrets,
credentials, and confidential documents. Such systems naturally
attract attackers, and we believe they are the most valuable to
defend. In contrast, disposable, single-use programs that are written
once and discarded fall outside our scope. This thesis focuses on
defending emerging software systems against the two untrusted factors
introduced above: AI-assisted coding tools at development time, and
autonomous agents at deployment time. We concretize this goal in two
domains: the security of Web3 software, and the mitigation of harmful
behavior in LLM-based systems.}
